\documentclass[11pt,a4paper]{article}

\usepackage[margin=25mm]{geometry}
\usepackage[T1]{fontenc}
\usepackage{lmodern}
\usepackage{amsmath,amssymb,amsthm}

\theoremstyle{definition}
\newtheorem{definition}{Definition}
\newtheorem{problem}{Problem}
\newcommand{\Area}{\operatorname{area}}
\setlength{\parindent}{0pt}
\setlength{\parskip}{6pt}

\begin{document}
\begin{center}
  {\Large\bfseries Heilbronn's Triangle Problem: Ten Points}
\end{center}
\section{Definition}
Let $U=[0,1]^2$. For a finite set $P\subseteq U$, define
\begin{equation}\label{eq:delta}
 h(P)
 =\min_{\substack{x,y,z\in P\\\text{distinct}}}
   \Area\bigl(\Delta(x,y,z)\bigr),
\end{equation}
where $\Delta(x,y,z)$ denotes the triangle formed by the three points $x,y,z$,
and degenerate triangles have area zero.

\begin{definition}[The $n$th Heilbronn number]
For $n\geq 3$, the \emph{$n$th Heilbronn number} of the unit square is
\begin{equation}\label{eq:H}
 H(n)=\max_{\substack{P\subseteq U\\|P|=n}} h(P).
\end{equation}
A set attaining this maximum is called an \emph{optimal Heilbronn
configuration}.
\end{definition}

The maximum exists by compactness. The use of the unit square is convenient but
not essential for asymptotic questions: for any two fixed convex bodies with
nonempty interior, the corresponding Heilbronn numbers differ by multiplicative
constants depending only on the two bodies.


\section{Current best result}
Optimal configurations have been determined or globally certified for
$3\leq n\leq 9$. The first unresolved case is $n=10$. 
The best known configurations satisfy
\begin{align*}
    H(10)&\geq 0.04653741958254,\\
H(11)&\geq \frac{1}{27},\\
H(12)&\geq 0.03259885869182.
\end{align*}

None of these values has been proved globally optimal.

\section{problem formalization}
Determine the exact value of $H(10)$ and classify all ten-point configurations
in the unit square that attain it.

\end{document}
